%% =================================================================
%% Wei, Mei, Zhao, Xiao, Sun, Zhang, Guo.
%% "Nondestructively identifying the mechanical behavior of soft
%%  tissues using surface deformation with an explicit inverse
%%  approach."
%% Appl. Math. Modelling 134 (2024) 126--147.
%%
%% Reconstruction of page 6 (printed p. 131) as study notes.
%% Prose: extracted verbatim from the PDF text layer via pdftotext.
%% Math:  hand-transcribed from the rendered page image — Eqs. (16)–(17)
%%        (adjoint equation + gradient of the objective function),
%%        plus two inline formulas (noise level, average relative error).
%% Figure 5: placeholder box describing the L-BFGS/adjoint flowchart.
%% =================================================================

\documentclass[11pt]{article}

\usepackage[margin=1in]{geometry}
\usepackage{amsmath, amssymb}
\usepackage{mathrsfs}
\usepackage{bm}
\usepackage{fancyhdr}

\pagestyle{fancy}
\fancyhf{}
\lhead{\small Y.~Mei et al.}
\rhead{\small Applied Mathematical Modelling 134 (2024) 126--147}
\cfoot{\small \thepage}
\renewcommand{\headrulewidth}{0.4pt}

\setcounter{page}{131}
\setcounter{equation}{15}  % next equation is (16)

\begin{document}

\noindent
The optimization problem is solved by the limited
Broyden--Fletcher--Goldfarb--Shanno (L-BFGS) method which requires
the objective function and the gradients of the objective function
with respect to the optimization variables. The gradients of the
objective function with respect to the optimization variables are
obtained by the adjoint method which has been thoroughly discussed
in [27]. The flow chart of the procedure to solve the inverse
problem is shown in Fig.~5.

To test the feasibility of the proposed inverse scheme, we adopt the
displacement datasets generated by a forward finite element
simulation with a target shear modulus distribution. The white
Gaussian random is also added into the simulated surface displacement
field. The noise level is defined by
$\sqrt{\displaystyle\sum_{k=1}^{N_{\text{node}}}
\left(u_k - u_k^{\text{noise}}\right)^{\!2}\bigg/
\sum_{k=1}^{N_{\text{node}}} (u_k)^{2}} \times 100\,\%$,
where $u_k$ and $u_k^{\text{noise}}$ are the simulated and noisy
displacement of $k$-th node, $N_{\text{node}}$ is the total number of
nodes in interest domain. In comparison, we also solve inverse
problems using the prevalent implicit inverse method based on the
adjoint method where the shear moduli are nodally defined [22]. To
quantitatively analyze the performance of the proposed method, we
define the average relative error by
$\dfrac{1}{N_{\text{node}}}
\displaystyle\sum_{k=1}^{N_{\text{node}}}
\left|\left(\mu_k - \mu_k^{\text{target}}\right)
\big/ \mu_k^{\text{target}}\right| \times 100\,\%$\,,
where $\mu_k$ is the recovered shear modulus value of $k$-th node
and $\mu_k^{\text{target}}$ is the target shear modulus value of
$k$-th node.

\subsection*{\textit{Gradient computation}}

\noindent
To efficiently compute the gradient of the objective function with
respect to the optimization variable, we have adopted the adjoint
method, which involves solving the associated adjoint equation. The
adjoint equation is given by [34,35]:
\begin{equation}
\mathscr{B}\!\left(\mathbf{w}^h, q^h, \delta\mathbf{u}^h, \delta p^h
;\, \mathbf{u}^h, p^h\right)
= -\bigl(\delta\mathbf{u}^h,\,
         \mathbf{u}^h - \mathbf{u}^{m}\bigr)_{\partial\Omega},
\qquad
\forall\,[\delta\mathbf{u}^h, \delta p^h]
\in \mathscr{M}^h \times \mathscr{P}^h
\end{equation}
where
$\mathscr{B}(\mathbf{w}^h, q^h, \delta\mathbf{u}^h, \delta p^h;\,
\mathbf{u}^h, p^h) = \frac{d}{d\varepsilon}
\mathscr{A}^{*}\!(\mathbf{w}^h, q^h, \mathbf{u}^h
+ \varepsilon\,\delta\mathbf{u}^h, p^h
+ \varepsilon\,\delta p^h)\big|_{\varepsilon=0}$
is the linearization of $\mathscr{A}^{*}$. With the acquired state
variables $\mathbf{u}^h, p^h$, we can solve for the dual variables
$\mathbf{w}^h, q^h$ using Eq.~(16). To this end, the gradient of the
objective function with respect to the shear modulus:
\begin{equation}
\delta\Pi = D_{\mu}\Pi\,\delta\mu = \mathscr{C}
+ \alpha \int_{\Omega}
\dfrac{\nabla\mu \cdot \nabla\delta\mu}
      {\sqrt{|\nabla\mu|^{2} + c^{2}}}\, d\Omega
\end{equation}
$\mathscr{C} = \frac{d}{d\varepsilon}
\mathscr{A}^{*}\!(\mathbf{w}^h, q^h, \mathbf{u}^h, p^h,
\mu + \varepsilon\,\delta\mu)\big|_{\varepsilon=0}$.
This can be used to efficiently compute the gradient of the objective
function with respect to

% ---------- Figure 5 placeholder ----------
\vspace{1em}
\begin{center}
\fbox{\parbox{0.92\textwidth}{\centering\vspace{1em}
\textbf{[Figure 5 — placeholder: optimization flowchart]}\\[0.8em]
\begin{tabular}{l}
\fbox{Give initial guess: $\mathbf{r}^1, \mathbf{r}^2, \ldots,
      \mathbf{r}^m$ and $\mu^0, \mu^1, \ldots, \mu^m$} \\[0.6em]
\hspace{2em}$\downarrow$\quad \textit{Finite Element Method} \\[0.4em]
\fbox{Solve forward problem $\to$ computed displacement
      $\mathbf{u}^{\text{com}}$} \\[0.6em]
\hspace{2em}$\downarrow$\quad \textit{Adjoint Method} \\[0.4em]
\fbox{Evaluate $\Pi$ and its gradients w.r.t.\ optimization
      variables} \\[0.6em]
\hspace{2em}$\downarrow$ \\[0.4em]
\fbox{Is $\Pi$ minimized?}
$\xrightarrow{\textit{No}}$
\fbox{Update variables (\textit{L-BFGS})}
$\to$ loop back $\uparrow$ \\[0.6em]
\hspace{2em}$\downarrow$\quad \textit{Yes} \\[0.4em]
\fbox{Stop and output reconstruction results}
\end{tabular}
\vspace{1em}}}
\end{center}
\vspace{0.3em}
\begin{center}
\textbf{Fig. 5.} The flow chart of the proposed explicit inverse method.
\end{center}

\end{document}
